% Source: https://artofproblemsolving.com/community/c3h2507477_fun_challenge_problem
% Stable identity: aops-topic-2507477-q01
A circle centred at $O$ with radius $r$ is not crossed by the line exactly when $r$ is at most the distance $d$ from $O$ to the line (for $r=d$ the line is tangent). So the largest area is $\pi d^2$ (in the sketch, $d=OF$).

To find $d$, drop the perpendicular from $O$ to the line. The line has direction $(1,m)$, so the perpendicular has direction $(m,-1)$, and its foot is $F=t(m,-1)$ for the $t$ with $-t=m\cdot tm+b$, i.e. $t=-\dfrac{b}{m^2+1}$. Hence
\[d=|t|\sqrt{m^2+1}=\dfrac{|b|}{\sqrt{m^2+1}},\qquad\text{and the largest area is}\qquad\boxed{\dfrac{\pi b^2}{m^2+1}}.\]

\begin{asy}
usepackage("amsmath");
size(5.5cm);
defaultpen(fontsize(10pt));
draw((-2.2,0)--(3.6,0),Arrow(TeXHead));draw((0,-2)--(0,3.1),Arrow(TeXHead));
draw(circle((0,0),1.6));
draw((-1,2.75)--(3.5,-0.625),linewidth(0.8));
pair F=(0.96,1.28);
draw((0,0)--F,dashed);dot(F);dot((0,2));
pair u=unit(-F),v=unit((4,-3));
draw(F+0.18*u--F+0.18*u+0.18*v--F+0.18*v);
label("$(0,b)$",(0,2),W,fontsize(9pt));
label("$d$",F/2,NW);
label("$y=mx+b$",(3.5,-0.625),E,fontsize(9pt));
label("$O$",(0,0),SW);
\end{asy}
